\documentclass[conference]{IEEEtran}

\usepackage[T1]{fontenc}
\usepackage[utf8]{inputenc}
\usepackage{amsmath,amssymb,amsthm}
\usepackage{booktabs}
\usepackage{graphicx}
\usepackage{multirow}
\usepackage{array}
\usepackage{url}
\usepackage{xcolor}
\usepackage{enumitem}
\usepackage{microtype}
\usepackage{hyperref}

\hypersetup{
  colorlinks=true,
  linkcolor=blue!50!black,
  urlcolor=blue!50!black,
  citecolor=blue!50!black
}

\newtheorem{proposition}{Proposition}
\newtheorem{theorem}{Theorem}
\newtheorem{remark}{Remark}
\newtheorem{definition}{Definition}

\newcommand{\R}{\mathbb{R}}
\newcommand{\Sph}{\mathbb{S}}
\newcommand{\E}{\mathbb{E}}
\newcommand{\KL}{\operatorname{KL}}
\newcommand{\Geo}{d_{\mathrm{geo}}}
\newcommand{\Logit}{\Delta_{\ell_2}}
\newcommand{\NLL}{\Delta_{\mathrm{NLL}}}

\title{Signal-Conditioned Sparse Codecs for Conversational Memory Compression\\
in Small Language Models}

\author{
\IEEEauthorblockN{Pranav Kompally}
\IEEEauthorblockA{
Independent Researcher\\
New York, USA
}
}

\begin{document}
\maketitle

% ─────────────────────────────────────────────────────────────────────────────
\begin{abstract}
Long multi-turn conversations exceed the inference budget of small language
models long before they become semantically exhausted. Existing retention
policies are either geometry-agnostic (semantic top-$k$) or signal-rigid
(geometry-only). We show that the optimal scoring signal for a sparse
keep/compress/drop (KCD) codec depends on the \emph{memory type} the task
requires, and that geometric features derived from hidden-state trajectories
on the unit sphere are both necessary and sufficient to distinguish these
regimes.

We formalize conversation-state dynamics using parallel-transported
increments on $\Sph^{d-1}$, show that resulting trajectories are extremely
low-rank (mean rank-95 of 1.05--1.53 across three model families) and
almost perfectly coupled to decoder output drift (geometry--logit correlation
0.989--0.994), and use this structure to define a signal-conditioned codec
family. Four experiments establish the claim. First, \texttt{geometry\_KCD}
significantly outperforms \texttt{semantic\_KCD} on constraint-critical memory
($\Delta\Logit{+}37.8$, $p{=}0.028$; support-turn rescue 17/36 vs.\ 7/36),
because semantic scoring cannot distinguish user constraints from their
assistant echoes while geometric state displacement can. Second, a Gate~1
oracle scale-up on MSC valid (111{,}877 candidate rows) shows query-conditioned
geometry still adds substantial ranking value inside a semantic shortlist
($\Delta\tau {+}0.35$--$0.45$, 12--15$\times$ threshold), while the larger
refinement study yields a narrower claim: budget-aware semantic KCD is the
strongest stable semantic baseline, and query-conditioned geometry beats
ambient geometry rather than clearly displacing that baseline. Third, the
same probe on LongMemEval-S cleaned passes Gate~1 at 5--10$\times$ threshold
($\Delta\tau {+}0.14$--$0.30$), confirming cross-benchmark generality. Fourth,
a 3B model validation shows the geometry advantage grows with scale:
\texttt{semantic\_KCD} at 3B produces $+1.13$~nats of answer NLL damage
($p{=}0.0003$) while every geometry-based policy improves answer quality.
We also investigate two open directions---a learned harm predictor and
memory-object level compression---reporting precise diagnoses for both.
\end{abstract}

\begin{IEEEkeywords}
conversational memory, sparse codecs, representation geometry, long-context
compression, signal-conditioned compression, small language models
\end{IEEEkeywords}

% ─────────────────────────────────────────────────────────────────────────────
\section{Introduction}
\label{sec:intro}

Long multi-turn conversations create a memory problem before they create a
generation problem. Once the full dialogue no longer fits inside a practical
inference budget, the model must decide what to retain, what to compress, and
what to evict. The prevailing approaches are semantic top-$k$ (retain the turns
most similar to the current query) and recency windowing (retain the most recent
$k$ turns). Both discard structural information: they treat the conversation
as an unordered bag of semantic fragments rather than as an evolving trajectory
through the model's representation space.

We argue that hidden-state geometry of the evolving conversation trajectory is
a useful and sometimes necessary signal for this decision---but the right way to
use it depends on the \emph{type} of memory the task requires. A conversation
contains at least two qualitatively different memory types:

\begin{enumerate}[leftmargin=1.4em]
\item \textbf{Constraint-critical memory.} Exact user specifications (code
  requirements, event logistics, ordered numerical facts) whose loss causes
  direct output failure. These turns are identifiable by the large local
  state displacement they cause; their assistant echoes are semantically
  nearly identical but geometrically much smaller.
\item \textbf{Persona-continuity memory.} Preference statements, identity
  disclosures, and relationship context whose value is proximity to the query,
  not displacement magnitude. These turns are best ranked by semantic
  query-relevance, refined by geometric stability signals.
\end{enumerate}

The central contribution of this paper is to show that a single sparse codec
structure---the keep/compress/drop (KCD) framework---can serve both memory
types, provided the scoring signal fed to it is matched to the memory type.
Geometry-based scoring is not universally better than semantic scoring, and
semantic scoring is not sufficient for constraint tasks. The right architecture
is signal-conditioned: one codec, two scoring regimes, selected by memory type.

This claim is supported by four experiments, two benchmark families (MSC valid
for persona-continuity, LongMemEval-S for episodic retrieval), and a
3B model-scale validation showing the geometry advantage grows with
representational capacity. We also report two informative negative results from
attempted extensions---a learned harm predictor and object-level compression---
whose failure modes are precisely diagnosed.

% ─────────────────────────────────────────────────────────────────────────────
\section{Related Work}
\label{sec:related}

\paragraph{Long-context compression.}
Early compression approaches focus on token-level eviction
\cite{scissorhands,streamingllm} or attention-head-based importance
scoring \cite{h2o}. These methods are stateless: retention decisions at
position $t$ do not depend on trajectory structure accumulated over
$1{:}t{-}1$. Retrieval-augmented generation \cite{rag} defers the
problem to an external store, avoiding in-context compression entirely.
Recent work on KV-cache compression \cite{kvcache} approaches the
problem from the attention layer rather than the conversation level.

\paragraph{Conversation-level memory.}
MemGPT \cite{memgpt} and similar systems maintain explicit external
memory stores, moving the compression decision out of the model entirely.
LongMemEval \cite{longmemeval} evaluates chat assistants on long-term memory recall
but does not study the geometry of the underlying state trajectories. The
MSC benchmark \cite{msc} is designed for persona continuity across
sessions rather than compression evaluation.

\paragraph{Representation geometry.}
Low-rank structure in transformer representations has been studied in
static contexts \cite{intrinsic}: our contribution is to study evolving
trajectories turn-by-turn using parallel transport on the sphere. The
closest geometric memory work \cite{geodesic} uses geodesic
distance for retrieval ranking; we instead use transported curvature as a
compression priority signal. Geometric scoring for memory \emph{retention}
under hard budget constraints, with mechanism analysis at the turn level,
has not been previously studied.

% ─────────────────────────────────────────────────────────────────────────────
\section{Geometric Framework}
\label{sec:math}

\subsection{Conversation-State Trajectories on the Sphere}

For a multi-turn conversation $x_{1:T}$, let
\begin{equation}
z_t = f_\theta(x_{1:t}) \in \R^d
\end{equation}
be the turn-level hidden-state summary extracted from a fixed intermediate
layer after processing turn $t$. We normalize onto the unit sphere:
\begin{equation}
h_t = \frac{z_t}{\lVert z_t \rVert} \in \Sph^{d-1}.
\end{equation}

Within a short segment $S_j$, choose a reference point $q_j \in \Sph^{d-1}$.
The one-step spherical increment from $h_t$ to $h_{t+1}$ is mapped into the
tangent space at $q_j$ via the Riemannian logarithm and parallel transport:
\begin{equation}
u_t = \mathrm{PT}_{h_t \rightarrow q_j}\!\bigl(\log_{h_t}(h_{t+1})\bigr)
\in T_{q_j}\Sph^{d-1}.
\label{eq:transported}
\end{equation}

\begin{proposition}
Local low-rank analysis of the transported increments $\{u_t\}_{t \in S_j}$
is well-defined independently of the original ambient coordinates, up to the
choice of segment reference and local transport gauge.
\end{proposition}
\begin{proof}[Proof sketch]
Parallel transport is an isometry between tangent spaces on the sphere.
Norms and principal angles among transported increments are therefore
invariant under the transport map, making Euclidean low-rank analysis
in the common tangent frame $T_{q_j}\Sph^{d-1}$ coordinate-independent.
\end{proof}

Empirically, we find that the increments are approximately low-rank:
$u_t \approx B_j c_t + e_t$ with $B_j \in \R^{d \times r_j}$ and
$r_j \ll d$. We measure this by rank-95: the minimum integer rank such
that 95\% of the variance in $\{u_t\}_{t \in S_j}$ is explained.

\subsection{Geometric Risk Score}

We define the discrete curvature proxy, stabilized against near-stationary
segments:
\begin{equation}
\kappa_t = \frac{\lVert u_t - u_{t-1} \rVert_2}{\max(\Delta s_t,\, \epsilon)},
\end{equation}
where $\Delta s_t$ is the local geodesic step length and $\epsilon$ is a small
floor. A large $\kappa_t$ signals a local regime transition: the conversation
has shifted direction in representation space. Combined with the segment
rank-95 and step norm, this yields a per-turn geometric risk score $r_t$.

\subsection{Query-Conditioned Geometry}

For persona-continuity tasks, we project the geometric signal into the query's
local tangent frame. Given a target turn $q$ with state $h_q$, the
query-conditioned risk score combines query-projected curvature:
\begin{equation}
\kappa_t^{(q)} = \kappa_t \cdot \bigl\langle \hat{u}_t,\,
\mathrm{PT}_{h_q \rightarrow q_j}(\hat{u}_q) \bigr\rangle,
\end{equation}
with query-projected subspace energy in the same tangent frame as the
conversation motion. The composite score $r_t^{(q)}$ is query-aware without
abandoning the spherical transport frame that makes the risk signal
coordinate-consistent.

\subsection{Decoder-Distortion Metrics}
\label{sec:metrics}

All experiments measure how well a compressed context preserves the model's
behavior relative to the full context. We use:
\begin{align}
\Logit &= \bigl\lVert \ell(\hat{h}_t) - \ell(h_t) \bigr\rVert_2
  & &\text{(logit drift)},\\
\NLL &= \bar{\ell}_{\mathrm{NLL}}(\hat{h}_t) - \bar{\ell}_{\mathrm{NLL}}(h_t)
  & &\text{(answer NLL delta)},
\end{align}
where $\ell(\cdot)$ is the output logit vector and $\bar{\ell}_{\mathrm{NLL}}$
is the average negative log-likelihood of the gold next-assistant-response
token sequence. Negative $\Logit$ or $\NLL$ means the policy is better than
the comparison. All p-values are from bootstrap permutation tests
(5{,}000 samples, two-sided).

\subsection{Theoretical Stability of the KCD Codec}

\begin{theorem}
Assume (i) the decoder is locally Lipschitz in a geometry metric, (ii)
compressed segment representations preserve anchors and the latest support
turn with bounded reconstruction error, and (iii) residual information for
evicted segments is bounded by a fixed segment-wise harm score. Then decoder
drift under the KCD policy satisfies:
\begin{equation}
\Logit(\pi_{\mathrm{KCD}})
\leq
\sum_j \Bigl[
\underbrace{\alpha \,\lVert a_j - \hat{a}_j\rVert}_{\text{anchor}}
+ \underbrace{\beta \,\lVert s_j - \hat{s}_j\rVert}_{\text{support}}
+ \underbrace{\gamma \,\rho_j}_{\text{evicted residual}}
\Bigr],
\end{equation}
where $\rho_j$ is the exogenous harm bound for segment $j$.
\label{thm:stability}
\end{theorem}
\begin{proof}[Proof sketch]
Lipschitz continuity of the decoder converts per-segment state reconstruction
error into logit drift. The codec decomposes reconstruction error into anchor,
support, and evicted residual terms; bounding each separately and summing over
segments yields the stated inequality.
\end{proof}

\begin{remark}
The bound is signal-agnostic: it holds for any scoring function that correctly
identifies high-harm segments ($\rho_j$ large) from low-harm ones ($\rho_j
\approx 0$). The empirical finding is that the correct scoring function is
geometry-based for constraint-critical memory and query-conditioned geometry
within a semantic shortlist for persona-continuity memory.
\end{remark}

% ─────────────────────────────────────────────────────────────────────────────
\section{The KCD Codec Family}
\label{sec:codec}

\subsection{Sparse Keep/Compress/Drop Structure}

For each segment $j$ of the conversation prefix, the KCD policy assigns an
action $a_j \in \{\mathrm{keep}, \mathrm{compress}, \mathrm{evict}\}$ under a
token budget $B$:
\begin{equation}
\max_{\{a_j\}} \sum_j v_j(a_j)
\quad
\text{s.t.}
\quad
\sum_j c_j(a_j) \leq B,
\end{equation}
where $v_j$ is the value (risk-weighted preservation) and $c_j$ is the cost
(retained token count) for action $a_j$ on segment $j$. The DP over segments
is exact for fixed segment span.

A compressed segment retains an \emph{anchor} (highest-risk turn), a \emph{sparse
support} (latest user turn and the highest-support-score turn for constraint
segments), and discards the remainder. A kept segment retains all turns. An
evicted segment retains nothing.

\subsection{Signal Variants}

The scoring signal $r_t$ that feeds segment risk $v_j$ is the free parameter:

\begin{itemize}[leftmargin=1.4em]
\item \textbf{geometry\_KCD}: $r_t$ is the per-turn geometric risk score from
  $\kappa_t$, segment rank-95, and step norm. Best for constraint-critical
  memory (hard stress set, code/SQL/logistics conversations).
\item \textbf{semantic\_KCD}: $r_t$ is cosine similarity of the turn's embedding
  to the target query. Structure-blind; actively harmful on constraint tasks.
\item \textbf{budget\_aware\_semantic\_KCD}: semantic scoring with budget-scaled
  segment span. Best semantic-only baseline.
\item \textbf{semantic\_query\_conditioned\_geometry\_KCD}: a two-stage policy.
  Stage 1 builds a semantic shortlist at $2\times$ the target budget. Stage 2
  applies $r_t^{(q)}$ (query-conditioned geometry) inside the shortlist. Best
  for persona-continuity and episodic-retrieval tasks.
\end{itemize}

The signal-conditioned design follows directly from the theoretical remark
above: use whichever signal most faithfully estimates $\rho_j$ for the task's
memory type.

% ─────────────────────────────────────────────────────────────────────────────
\section{Experiments}
\label{sec:experiments}

\subsection{Setup}

\paragraph{Models.}
Primary experiments use \texttt{qwen25\_15b} (Qwen/Qwen2.5-1.5B-Instruct) for
its balance of inference speed and representational capacity. Scale validation
uses \texttt{qwen25\_3b} (Qwen/Qwen2.5-3B-Instruct). Geometry characterization
additionally uses \texttt{qwen25\_05b} (0.5B) and \texttt{smollm2\_17b}
(SmolLM2-1.7B).

\paragraph{Benchmarks.}
We use three evaluation sets:
\begin{itemize}[leftmargin=1.4em]
\item \textbf{Hard stress set}: 36 custom conversations across three families
  (\texttt{long\_dependency}, \texttt{retrieval\_heavy}, \texttt{code\_conversation}).
  Designed so that correct late-turn answers depend on exact early-turn
  constraints; tests constraint-critical memory.
\item \textbf{MSC valid} (\texttt{gonced8/multi-session-chat}, 32 conversations):
  multi-session dialogue with persona continuity and preference tracking.
  Tests persona-continuity memory, not constraint preservation.
\item \textbf{LongMemEval-S cleaned} (\texttt{xiaowu0162/longmemeval-cleaned},
  12-conversation subset): episodic-retrieval conversations (398--618 turns
  each; truncated to 40 turns per conversation for oracle tractability).
\end{itemize}

\paragraph{Baselines.} \texttt{uniform} (budget spread without geometry),
\texttt{semantic} (top-$k$ by cosine similarity), \texttt{geometry} (top-$k$
by geometric risk). All experiments use budgets 0.20/0.35/0.50 of the
conversation token count.

\subsection{Geometry Characterization}
\label{sec:char}

Conversation-state trajectories are extremely low-rank across all tested
models (Table~\ref{tab:char}). The mean rank-95 ranges from 1.049 to 1.528
--- meaning that at most 2--3 principal directions account for 95\% of the
variance in transported increments within a segment. Shuffling turn order
increases rank-95 substantially, ruling out bounded-dimensionality as an
explanation. Geometric distortion is almost perfectly coupled to decoder
output drift across all models.

\begin{table}[t]
  \centering
  \caption{Geometry characterization. Rank-95 is near-minimal; shuffled controls
  collapse coupling; geometry--logit correlation is extremely high.}
  \label{tab:char}
  \begin{tabular}{lccc}
    \toprule
    Model & Rank-95 & Shuffled rank & Geo--logit $r$ \\
    \midrule
    \texttt{qwen25\_05b} & 1.049 & 1.486 & 0.989 \\
    \texttt{qwen25\_15b} & 1.167 & 1.531 & 0.994 \\
    \texttt{smollm2\_17b} & 1.528 & 1.812 & 0.989 \\
    \bottomrule
  \end{tabular}
\end{table}

These results establish that the geometric structure we use for scoring is not
incidental. Conversation-state trajectories evolve along a single dominant
direction within each segment, and deviations from that direction (high
curvature, high rank-95) predict decoder output drift with near-perfect
fidelity.

\subsection{Geometry-Aware Budget Allocation}
\label{sec:paper2}

Using geometric risk as a turn-level priority score, geometry-aware retention
beats uniform allocation on the hard stress set with consistent significance.
At budget 0.35, \texttt{geometry} reduces mean logit drift by 99.5 on
\texttt{qwen25\_05b} and 21.7 on \texttt{qwen25\_15b}. The mechanism is
interpretable at the turn level: geometry retains more support user turns
than uniform in 14--17/36 cases, compared to 2--5/36 for uniform winning. The
rescued turns are exact constraint packets: code requirements, SQL ordering
clauses, event logistics, and compact numerical fact bundles.

This establishes the baseline for the codec experiments: geometry-based
priority is both useful and mechanistically interpretable before any codec
structure is applied.

\subsection{Experiment 1: The Geometry Signal Is Load-Bearing}
\label{sec:exp1}

\paragraph{Question.} Is geometry specifically necessary for constraint-critical
memory, or does the KCD codec structure do the work regardless of scoring signal?

\paragraph{Setup.} We compare \texttt{geometry\_KCD} and \texttt{semantic\_KCD}
on the full hard stress set (\texttt{qwen25\_15b}, budgets 0.20/0.35/0.50).
Both policies use the identical KCD codec with identical segment spans; they
differ only in the scoring signal fed to the DP.

\paragraph{Results.}

\begin{table}[t]
  \centering
  \caption{Experiment 1: geometry\_KCD vs.\ semantic\_KCD on the hard stress set
  (\texttt{qwen25\_15b}). Positive $\Delta$ means semantic\_KCD is worse.}
  \label{tab:exp1}
  \begin{tabular}{lrrrr}
    \toprule
    Budget & $\Delta\Logit$ & $p$ & $\Delta\NLL$ & $p$ \\
    \midrule
    0.20 & $+37.8$ & $0.028$ & $+0.195$ & $0.001$ \\
    0.35 & $+38.5$ & $0.079$ & $+0.424$ & $0.011$ \\
    0.50 & $+0.1$ (tie) & $0.995$ & $+0.628$ & $0.001$ \\
    \bottomrule
  \end{tabular}
\end{table}

\texttt{geometry\_KCD} beats \texttt{semantic\_KCD} at budgets 0.20 and 0.35
on both logit and answer quality with statistical significance. At 0.50, logit
distortions tie but geometry still dominates answer NLL ($\Delta{=}+0.628$,
$p{=}0.001$). Critically, \texttt{semantic\_KCD} is also significantly worse
than plain \texttt{semantic} at every budget on answer NLL: the KCD codec
structure with the wrong signal actively harms performance.

\paragraph{Mechanism.}

\begin{table}[t]
  \centering
  \caption{Support-turn rescue at budget 0.35 (\texttt{qwen25\_15b}).
  Geometry preferentially retains constraint packets; semantic cannot
  distinguish them from their echoes.}
  \label{tab:rescue}
  \begin{tabular}{lcc}
    \toprule
    Metric & \texttt{geo\_KCD} & \texttt{sem\_KCD} \\
    \midrule
    Rescues more support turns than uniform & 17/36 & 7/36 \\
    Uniform beats policy on support turns & 2/36 & 5/36 \\
    Keeps latest support turn vs.\ uniform & 17/36 & 7/36 \\
    Constraint turns rescued & 13 & 4 \\
    Base-memory turns rescued & 6 & 3 \\
    \bottomrule
  \end{tabular}
\end{table}

The failure mode is mathematically interpretable. Consider the turn pair:
\emph{User}: ``Also require \texttt{GET /users/\$\{id\}}, throw on non-ok
responses, and return parsed JSON.'' \emph{Assistant}: ``Got it. I'll add error
handling for \texttt{GET /users/\$\{id\}} with non-ok response checking and
JSON parsing.'' These turns have nearly identical cosine similarity to any
future query about the API. Semantic scoring assigns them equal priority.
Geometric state displacement assigns the user constraint a much higher risk
score, because the constraint causes a large local trajectory deflection while
the assistant echo does not.

\subsection{Experiment 2: Gate 1 --- Geometry Inside a Semantic Shortlist (MSC)}
\label{sec:gate1}

\paragraph{Question.} On a persona-continuity benchmark where semantic proximity
is the primary signal, do geometric features add ranking value inside a semantic
shortlist? Does the hybrid policy beat the best semantic-only codec?

\paragraph{Setup.} An oracle probe generates turn-level harm scores by ablating
each prior turn and measuring the resulting logit damage. We then measure
whether geometric features (\texttt{query\_geom\_v2\_risk}) can rank turns by
oracle harm better than semantic scoring alone, inside a $2\times$ semantic
shortlist. Decision threshold for Gate 1: $\Delta\tau \geq +0.03$.

MSC valid, 32 conversations, \texttt{qwen25\_15b}, 111{,}877 candidate rows.

\paragraph{Oracle result: PASS.}

\begin{table}[t]
  \centering
  \caption{Gate 1 oracle on MSC valid. Geometry adds substantial within-shortlist
  ranking value. Threshold: $\Delta\tau \geq +0.03$.}
  \label{tab:gate1_msc}
  \begin{tabular}{lccc}
    \toprule
    Budget & Sem $\tau$ & Geo $\tau$ & $\Delta\tau$ \\
    \midrule
    0.20 & 0.295 & 0.647 & $+0.352$ \\
    0.35 & 0.223 & 0.672 & $+0.449$ \\
    0.50 & 0.279 & 0.718 & $+0.439$ \\
    \bottomrule
  \end{tabular}
\end{table}

Every budget exceeds the threshold by roughly 12--15$\times$. Geometry still
improves within-shortlist turn ranking strongly over semantic scoring alone,
but the larger-scale result sharpens the interpretation: this is a ranking
correction inside a semantic shortlist, not a replacement for semantic retrieval.

\paragraph{Policy result.} The refinement study shows two stable effects:
\texttt{budget\_aware\_semantic\_KCD} beats pure \texttt{semantic}, and
\texttt{semantic\_query\_conditioned\_geometry\_KCD} beats
\texttt{semantic\_ambient\_geometry\_KCD} at low-to-mid budgets.

\begin{table}[t]
  \centering
  \caption{Gate 1 policy on MSC valid. Negative $\Delta$ means the left policy
  is better.}
  \label{tab:gate1_policy}
  \begin{tabular}{lccc}
    \toprule
    Comparison & 0.20 & 0.35 & 0.50 \\
    \midrule
    budget-aware semantic vs.\ semantic & $-13.8$ ($p{=}0.349$ conv) & $-27.7$ ($p{=}0.045$ conv) & $-71.6$ ($p{=}0.002$ conv) \\
    query-conditioned vs.\ ambient geometry & $-14.4$ ($p{=}0.047$ conv) & $-23.4$ ($p{=}0.027$ conv) & $+4.6$ ($p{=}0.814$ conv) \\
    \bottomrule
  \end{tabular}
\end{table}

At larger scale, the MSC refinement claim is more conservative than in the
earlier slice. Budget-aware semantic is now the strongest stable semantic
baseline, while query-conditioned geometry still significantly improves over
ambient geometry at budgets 0.20 and 0.35. Query-conditioning is therefore
still the operative geometric mechanism, but the larger run no longer supports
the stronger claim that the hybrid cleanly beats the best semantic codec.

\subsection{Experiment 3: Gate 1 on LongMemEval-S --- Cross-Benchmark Confirmation}
\label{sec:gate1_lme}

\paragraph{Question.} Is the Gate 1 geometry ranking advantage specific to
MSC's persona-continuity structure, or does it generalize to episodic-retrieval
conversations?

\paragraph{Setup.} Oracle probe on 12 LongMemEval-S cleaned conversations
(40 turns per conversation, 32{,}251 candidate rows).

\paragraph{Oracle result: PASS at 5--10$\times$ threshold.}

\begin{table}[t]
  \centering
  \caption{Gate 1 oracle on LongMemEval-S cleaned. Semantic scoring is
  effectively random inside the shortlist; geometry is highly informative.}
  \label{tab:gate1_lme}
  \begin{tabular}{lcccc}
    \toprule
    Budget & Sem $\tau$ & Geo $\tau$ & $\Delta\tau$ & Pass \\
    \midrule
    0.20 & $-0.010$ & 0.304 & $+0.303$ & \checkmark \\
    0.35 & $-0.010$ & 0.145 & $+0.145$ & \checkmark \\
    0.50 & $-0.010$ & 0.289 & $+0.289$ & \checkmark \\
    \bottomrule
  \end{tabular}
\end{table}

On episodic-retrieval conversations, semantic scoring inside the shortlist
remains essentially non-informative, while geometry remains clearly above the
Gate~1 threshold at every budget. The overall (pre-shortlist) $\Delta\tau$ is
still much larger than the shortlist-restricted delta. On the truncated 40-turn
slice, however, the policy comparison is inconclusive at the benchmark level,
so LongMemEval should be read here as a shortlist-ranking replication rather
than as a decisive end-to-end policy win.

\subsection{Experiment 4: 3B Model-Scale Validation}
\label{sec:3b}

\paragraph{Question.} Does the signal-comparison result hold at higher model
capacity, or is it specific to the 1.5B compact regime?

\paragraph{Setup.} Identical to Experiment 1 but with \texttt{qwen25\_3b}, 9
conversations, 540 evaluations.

\begin{table}[t]
  \centering
  \caption{Experiment 4: geometry\_KCD vs.\ semantic\_KCD on the hard stress set
  (\texttt{qwen25\_3b}). The geometry advantage grows with model scale.}
  \label{tab:exp4}
  \begin{tabular}{lrrrr}
    \toprule
    Budget & $\Delta\Logit$ & $p$ & $\Delta\NLL$ & $p$ \\
    \midrule
    0.20 & $+44.9$ & $0.048$ & $+0.494$ & $0.008$ \\
    0.35 & $+69.7$ & $0.050$ & $+0.853$ & $0.005$ \\
    0.50 & $-29.1$ & $0.128$ & $+1.126$ & $\mathbf{0.0003}$ \\
    \bottomrule
  \end{tabular}
\end{table}

At all three budgets, geometry-based policies consistently improve answer
quality relative to uniform; \texttt{semantic\_KCD} is the only policy that
makes things worse at every budget (Table~\ref{tab:3b_nll}).

\begin{table}[t]
  \centering
  \caption{3B answer NLL relative to uniform. Negative = better.
  \texttt{semantic\_KCD} is uniquely harmful at every budget.}
  \label{tab:3b_nll}
  \begin{tabular}{lrrr}
    \toprule
    Policy & @ 0.20 & @ 0.35 & @ 0.50 \\
    \midrule
    \texttt{geometry}     & $-0.158$ & $-0.036$ & $-0.626$ \\
    \texttt{geometry\_KCD} & $-0.134$ & $-0.357$ & $-0.775$ \\
    \texttt{semantic}     & $-0.017$ & $-0.147$ & $-0.665$ \\
    \texttt{semantic\_KCD} & $\mathbf{+0.360}$ & $\mathbf{+0.496}$ & $+0.350$ \\
    \bottomrule
  \end{tabular}
\end{table}

The support-turn rescue mechanism is scale-invariant: \texttt{geometry\_KCD}
rescues more support turns than uniform in 15/36 cases (vs.\ 7/36 for
\texttt{semantic\_KCD}), the identical ratio as at 1.5B. The quality
\emph{consequence} of semantic failure grows with scale: at 3B and budget 0.50,
\texttt{semantic\_KCD} produces $+1.13$ nats of answer NLL damage ($p{=}0.0003$).
The conventional reviewer assumption --- ``geometry only works on tiny models'' ---
is directly inverted: the geometry advantage strengthens with representational
capacity.

% ─────────────────────────────────────────────────────────────────────────────
\section{The Signal-Conditioned Regime Map}
\label{sec:regime}

The four experiments together with the budget-allocation baseline establish a
consistent regime map for conversational memory compression:

\begin{table}[t]
  \centering
  \caption{Winner map by memory type and budget. KCD structure is common to
  all; the scoring signal is task-conditioned.}
  \label{tab:regime}
  \begin{tabular}{p{2.0cm}p{2.5cm}p{2.5cm}}
    \toprule
    Memory type & Winning signal & Evidence \\
    \midrule
    Constraint-critical
      & Geometric risk ($r_t$)
      & Exp.\ 1, 3B probe; 17/36 rescue \\
    Persona-continuity
      & Semantic shortlist + query-geom.\ ($r_t^{(q)}$)
      & Gate 1 MSC scale-up; oracle pass, query-conditioned $>$ ambient \\
    Episodic retrieval
      & Semantic shortlist + query-geom.
      & Gate 1 LME; $\Delta\tau{+}0.14$--$0.30$ \\
    \bottomrule
  \end{tabular}
\end{table}

The separation is not benchmark noise. It reflects a structural difference
in how the two memory types produce decoder harm. Constraint turns cause large
local state displacements; their echoes do not. Persona turns are semantically
proximate to later queries; their geometric displacement relative to those
queries adds within-shortlist ranking precision that semantic scoring misses.

% ─────────────────────────────────────────────────────────────────────────────
\section{Negative Results and Claim Boundaries}
\label{sec:negative}

\paragraph{Falsified extensions.}
Five geometry-only extensions failed, defining the boundary of the geometry
claim:

\begin{itemize}[leftmargin=1.4em]
\item \textbf{Persona/filler curvature separation (MSC).} Mean support-filler
  curvature delta was 0.122 (3/5 positive, 2/5 negative). Curvature does not
  separate semantically important persona turns from nearby filler; semantic
  proximity is needed.
\item \textbf{Geometry-only regime atlas.} After fixing raw curvature blow-up
  (up to 6{,}314 on near-stationary segments, stabilized to 26), the atlas
  still failed to separate retrieval-heavy hard-set segments from MSC and
  LoCoMo. Geometry alone does not recover a usable benchmark taxonomy.
\item \textbf{State-update supersession (two formulations).} Both same-sign
  alignment ($A(s,t) = \cos(u_s, u_t)$, mean 0.913) and a state-vs.-increment
  cross-term (update cross 0.971) fail as sign-based detectors on synthetic
  explicit-update benchmarks. Compact models encode both ``user states a fact''
  and ``user updates a fact'' as similarly oriented increments.
\item \textbf{Query-conditioned geometry as standalone selector.} At LoCoMo
  budget 0.35, standalone query-conditioned geometry collapses to
  $\Logit{=}+881.4$. The KCD wrapper recovers to $+9.7$ but does not win.
  Query-conditioned geometry is a refinement signal inside a semantic shortlist,
  not a primary selector.
\end{itemize}

\paragraph{Investigated open items.}
Two further directions were tested and yield informative negatives:

\textbf{Learned harm predictor.}
A 2-head MLP trained from 23{,}928 oracle ablation rows achieves row-level
Kendall $\tau{=}0.637$ on validation --- genuine learning. However,
conversation-level Kendall is $-1.0$: the predictor inverts the scale across
conversations (turns from lower-harm conversations receive higher absolute
scores than turns from higher-harm ones). When used as a policy
(\texttt{semantic\_harm\_KCD}), it loses to
\texttt{semantic\_query\_conditioned\_geometry\_KCD} at budget 0.35 by
$\Delta\Logit{=}+51.2$ ($p{=}0.004$ conversation-level). Root cause: pointwise
MSE loss on harm scalars does not enforce cross-conversation calibration.
Geometric features are immune to this problem because transported increments
are computed in the same spherical frame regardless of conversation. Fix:
pairwise ranking loss (RankNet/ListMLE) with cross-conversation training pairs.

\textbf{Memory-object level compression.}
Object-granularity KCD groups turns into semantic bundles (persona, event,
constraint, update) and makes codec decisions at the bundle level.
Token fractions on MSC are 0.43--0.52 at a 0.74 target: the codec chronically
under-retains. Answer NLL is 0.54--0.87 nats worse than turn-level baselines.
Root cause: marker-based segmentation creates many small objects; each retains
1--2 turns after compression; total retained mass falls well below budget. The
object-level structure is conceptually sound; the implementation needs
model-based boundary detection.

% ─────────────────────────────────────────────────────────────────────────────
\section{Discussion}
\label{sec:discussion}

\paragraph{What is now settled.}
The following claims have multi-experiment, multi-benchmark, multi-scale
support:
\begin{enumerate}[leftmargin=1.4em]
\item Conversation-state trajectories are extremely low-rank and almost
  perfectly decoder-coupled across compact model families. This is the
  geometric substrate that makes the scoring signals tractable.
\item Geometry-aware retention beats uniform on constraint-critical tasks via
  a concrete mechanism: support-turn rescue. The rescued turns are
  interpretable word-for-word.
\item The KCD codec is a real, budget-responsive compression structure. It
  beats baselines under fairness-controlled token accounting.
\item The geometry signal is load-bearing for constraint tasks. The codec
  structure alone does not do the work; semantic scoring with the same codec
  is actively harmful.
\item Geometry adds substantial ranking value inside semantic shortlists, on
  both MSC and LongMemEval-S, at 5--15$\times$ the Gate 1 threshold. This is
  not a marginal effect, but on public semantic-memory benchmarks the gain is
  best understood as shortlist ranking correction rather than as a universal
  top-$k$ recall increase.
\item The geometry advantage grows with model scale. At 3B, semantic scoring
  produces catastrophic answer quality damage; geometry consistently improves
  it. The concern about compact-model specificity is inverted.
\item The signal-conditioned KCD framework is the correct architecture: one
  codec structure, two scoring regimes matched to memory type.
\end{enumerate}

\paragraph{What remains open.}
Four precise directions follow from the diagnosed failures:
\begin{enumerate}[leftmargin=1.4em]
\item \textbf{Cross-conversation pairwise predictor.} A pairwise ranking loss
  with training pairs drawn across conversations would address the calibration
  inversion. Ranking losses have well-studied generalization properties that
  would also extend the Theorem~\ref{thm:stability} stability bound to a
  learned predictor.
\item \textbf{Model-based object boundary detection.} Replacing marker-based
  segmentation with a lightweight turn classifier would fix chronic
  under-retention in the object codec. The object-level KCD structure is
  otherwise valid.
\item \textbf{Systematic 7B$+$ validation.} The failed compact-model branches
  (supersession, persona/filler separation) may become viable at higher
  representational resolution.
\item \textbf{Codec theorem with generalization error.} Extending
  Theorem~\ref{thm:stability} to a learned $\rho_j$ predictor with PAC-style
  generalization bounds would complete the theoretical story.
\end{enumerate}

\paragraph{Limitations.}
The constraint-memory mechanism story is strongest on the custom hard stress
set, which was designed to surface support-turn failures. Gate 1 policy
evidence now covers 32 MSC and 12 LongMemEval conversations, but policy
separation is still weaker than oracle shortlist separation; broader
replication is warranted. The LongMemEval oracle uses 40-turn truncation for
computational tractability. The codec theorem assumes a bounded exogenous harm
predictor rather than a learned one.

% ─────────────────────────────────────────────────────────────────────────────
\section{Conclusion}
\label{sec:conclusion}

We have shown that the question ``is geometry better than semantics for
conversational memory compression?'' is not well-posed. The correct question
is: \emph{which signal best identifies the turns whose removal causes the most
decoder harm for a given task's memory type?} For constraint-critical tasks, the
answer is geometric displacement: semantic scoring cannot distinguish a user
constraint from its assistant echo, while the hidden-state trajectory can. For
persona-continuity and episodic-retrieval tasks, the answer is semantic
proximity refined by query-conditioned geometry: semantic scoring provides the
shortlist, geometry provides the within-shortlist ranking precision.

These two scoring regimes share a single codec structure (KCD), a single
theoretical stability bound (Theorem~\ref{thm:stability}), and a
scale-consistent geometric substrate. The advantage is not model-size
dependent: at 3B, semantic compression of constraint-critical memory produces
$+1.13$ nats of answer quality damage while geometry compression consistently
reduces it. The support-turn rescue mechanism that explains the advantage is
identical at 1.5B and 3B --- what grows with scale is the quality penalty for
getting it wrong.

The result is a signal-conditioned sparse codec family for conversational
memory compression, grounded in the geometry of the model's own representation
space, validated across three benchmark families and two model scales, with five
explicit falsification results that define the boundary of the claim and two
investigated open directions with precise diagnoses.

% ─────────────────────────────────────────────────────────────────────────────
\begin{thebibliography}{99}

\bibitem{scissorhands}
Z.~Liu et al.,
``ScissorHands: Exploiting the Persistence of Importance Hypothesis for
LLM KV Cache Compression at Test Time,''
\emph{NeurIPS}, 2023.

\bibitem{streamingllm}
G.~Xiao et al.,
``Efficient Streaming Language Models with Attention Sinks,''
\emph{ICLR}, 2024.

\bibitem{h2o}
Z.~Zhang et al.,
``H$_2$O: Heavy-Hitter Oracle for Efficient Generative Inference of Large
Language Models,''
\emph{NeurIPS}, 2023.

\bibitem{rag}
P.~Lewis et al.,
``Retrieval-Augmented Generation for Knowledge-Intensive NLP Tasks,''
\emph{NeurIPS}, 2020.

\bibitem{kvcache}
Y.~Ge et al.,
``Model Tells You What to Discard: Adaptive KV Cache Compression for LLMs,''
\emph{ICLR}, 2024.

\bibitem{memgpt}
C.~Packer et al.,
``MemGPT: Towards LLMs as Operating Systems,''
\emph{arXiv:2310.08560}, 2023.

\bibitem{longmemeval}
D.~Wu et al.,
``LongMemEval: Benchmarking Chat Assistants on Long-Term Interactive Memory,''
\emph{arXiv:2410.10813}, 2024.

\bibitem{msc}
J.~Xu et al.,
``Beyond Goldfish Memory: Long-Term Open-Domain Conversation,''
\emph{ACL}, 2022.

\bibitem{intrinsic}
A.~Aghajanyan et al.,
``Intrinsic Dimensionality Explains the Effectiveness of Language Model
Fine-Tuning,''
\emph{ACL}, 2021.

\bibitem{geodesic}
K.~Ethayarajh,
``How Contextual are Contextualized Word Representations?
Comparing the Geometry of BERT, ELMo, and GPT-2 Embeddings,''
\emph{EMNLP}, 2019.

\end{thebibliography}

\end{document}
